% --- Archon physics formalization source begin ---
% archon:physics
% archon:covers PhyXMiniProblems/problem_phyx_mini_0776.lean
% archon:source-report reports/phyx_mini/problem_phyx_mini_0776.source.json

\chapter{Physics problem phyx\_mini\_0776}
\label{ch:PhyXMiniProblems_problem_phyx_mini_0776}

\paragraph{Problem source.}
\#\# Physical scenario

The cone has mass \textbackslash{}(M\textbackslash{}) and altitude \textbackslash{}(h\textbackslash{}). The radius of its circular base is \textbackslash{}(R\textbackslash{}).

\#\# Question

Calculate the moment of inertia of a uniform solid cone about an axis through its center as shown in figure.

\#\# Answer choices

- A: A: \textbackslash{}(\textbackslash{}frac\{3\}\{8\} MR\textasciicircum{}2\textbackslash{})

- B: B: \textbackslash{}(\textbackslash{}frac\{3\}\{10\} MR\textasciicircum{}2\textbackslash{})

- C: C: \textbackslash{}(\textbackslash{}frac\{7\}\{10\} MR\textasciicircum{}2\textbackslash{})

- D: D: \textbackslash{}(\textbackslash{}frac\{5\}\{8\} MR\textasciicircum{}2\textbackslash{})

\#\# Figure caption (auxiliary; use the image as primary evidence)

The image depicts a three-dimensional cone. 

1. **Objects and Structures**:
   - The cone is oriented with the point facing downwards and the circular base at the top.
   - A vertical black line represents the "Axis" of the cone, extending from the center of the base vertically through the tip.

2. **Text and Labels**:
   - The base radius is labeled as \textbackslash{}( R \textbackslash{}) with a pink arrow extending from the center to the edge of the base.
   - The height of the cone is labeled as \textbackslash{}( h \textbackslash{}), marked with a vertical black line extending from the center of the base to the tip.

3. **Colors and Styles**:
   - The cone is shaded in a gradient of orange.
   - The radius arrow is pink, providing contrast against the cone’s color.

4. **Geometry and Dimensions**:
   - The relationships between the base, height, and axis are clearly illustrated with perpendicular lines and labels, indicating the standard geometric properties of a cone.

Overall, the image effectively illustrates the basic dimensions and structure of a right circular cone.

\#\# Dataset metadata

Rotational Motion; Physical Model Grounding Reasoning, Numerical Reasoning

\paragraph{Recorded answer/context.}
B: B: \textbackslash{}(\textbackslash{}frac\{3\}\{10\} MR\textasciicircum{}2\textbackslash{})

\paragraph{Figure/image path.}
/root/proposal\_for\_physic/science-mango/phyx\_mini\_run/phyx\_data/test\_image/776.png

\paragraph{Formalization target.}
create a compiling Lean file with sorry bodies at `PhyXMiniProblems/problem\_phyx\_mini\_0776.lean`. The Lean declarations must preserve the physical quantities, dimensions or dimensional roles, figure labels, governing-law hypotheses, and final relation expressed by this problem.
Use Mathlib/Physlib names found through LeanExplore where available. If a physics API is missing, introduce faithful local abstractions rather than scalar placeholder aliases.

\begin{theorem}[Physics formalization target]
\label{thm:physics:phyx_mini_0776:target}
\lean{PhyXMini0776.uniformSolidCone_momentOfInertia_about_symmetryAxis}
\uses{def:physics:phyx-mini-0776:phyxmini0776-momentofinertiaaboutsymmetryaxis, def:physics:phyx-mini-0776:phyxmini0776-isuniformsolidrightcircularcone}
The assigned autoformalize agent should translate this physics problem into Lean declarations in the covered file, with theorem and lemma proof bodies written as `by sorry`.
\end{theorem}
\begin{proof}
This is an autoformalization task, not a proof task. Produce statements that can later be proved by the physics prover without weakening the physical model.
\end{proof}

% --- Archon declaration topology begin ---
\section{Declaration topology}
\begin{definition}[symmetry Axis Index]
  \label{def:physics:phyx-mini-0776:phyxmini0776-symmetryaxisindex}
  \lean{PhyXMini0776.symmetryAxisIndex}
  The coordinate direction of the cone's axis in the body-fixed frame from the figure
\end{definition}
\begin{proof}
  This is the declared model definition of symmetry Axis Index.
\end{proof}
\begin{definition}[symmetry Axis Line]
  \label{def:physics:phyx-mini-0776:phyxmini0776-symmetryaxisline}
  \lean{PhyXMini0776.symmetryAxisLine}
  The line labelled Axis in the figure: the third coordinate axis
\end{definition}
\begin{proof}
  This is the declared model definition of symmetry Axis Line.
\end{proof}
\begin{definition}[solid Right Circular Cone Region]
  \label{def:physics:phyx-mini-0776:phyxmini0776-solidrightcircularconeregion}
  \lean{PhyXMini0776.solidRightCircularConeRegion}
  \uses{def:physics:phyx-mini-0776:phyxmini0776-symmetryaxisindex}
  The solid right circular cone shown in the figure, with its apex at the origin, its base in the plane x₂ = altitude, and its radius there equal to radius. Coordinates and the values of altitude and radius are measured in the same length unit
\end{definition}
\begin{proof}
  This is the declared model definition of solid Right Circular Cone Region.
\end{proof}
\begin{definition}[moment Of Inertia About Symmetry Axis]
  \label{def:physics:phyx-mini-0776:phyxmini0776-momentofinertiaaboutsymmetryaxis}
  \lean{PhyXMini0776.momentOfInertiaAboutSymmetryAxis}
  \uses{def:physics:phyx-mini-0776:phyxmini0776-symmetryaxisindex}
  The physical moment of inertia about the line labelled Axis in the figure. The dimension tag records the dimension mass * length * length
\end{definition}
\begin{proof}
  This is the declared model definition of moment Of Inertia About Symmetry Axis.
\end{proof}
\begin{definition}[Is Uniform Solid Right Circular Cone]
  \label{def:physics:phyx-mini-0776:phyxmini0776-isuniformsolidrightcircularcone}
  \lean{PhyXMini0776.IsUniformSolidRightCircularCone}
  \uses{def:physics:phyx-mini-0776:phyxmini0776-solidrightcircularconeregion}
  The governing model for a uniform solid cone. The rigid body's mass distribution is constant on solidRightCircularConeRegion and zero outside it. The coefficient is the uniform density obtained from the cone volume π R² h / 3. This predicate contains the setup and governing law only; it does not assume a formula for the moment of inertia
\end{definition}
\begin{proof}
  This is the declared model definition of Is Uniform Solid Right Circular Cone.
\end{proof}
% --- Archon declaration topology end ---
% --- Archon physics formalization source end ---
